\section{Hai nhánh đánh giá và món nợ giữa chúng}
\label{sec:ch12-hai-nhanh-danh-gia}

\index{đánh giá!hai nhánh|(}%
\index{đánh giá theo chức năng}%
Ba cách đánh giá mà mục~\ref{sec:ch01-danh-gia-giai-thich} đọc không phải ba
nhãn chất lượng cho cùng một đại lượng, và chỗ quan trọng nhất trong cách
Doshi-Velez và Kim phát biểu chúng là quan hệ vay mượn giữa chúng. Đánh giá
theo chức năng không cần người tham gia, nên nó rẻ và chạy được hàng loạt;
nhưng hai tác giả nói rõ rằng nó thích hợp nhất khi lớp mô hình đang xét đã
được kiểm chứng từ trước bằng một cách đánh giá đắt hơn, tức bằng một thí
nghiệm có người~\autocite{p14rigorous}. Điều kiện ấy là điều kiện mà
mục~\ref{sec:ch01-danh-gia-giai-thich} đã phát biểu ở dạng tổng quát hơn: một
\tn{đại lượng thay thế}{proxy} chỉ đáng tin khi ai đó đã cho thấy nó đi cùng
điều người dùng thật sự cần.
Chương~\ref{ch:do-do-trung-thuc} đã xếp toàn bộ các họ \tn{chỉ số trung thực}{faithfulness metric} vào ô
thứ ba ấy.

Món nợ vừa nêu là điều chương này đi đòi. Nếu mọi chỉ số trung thực đều là
đại lượng thay thế thuộc cách đánh giá rẻ nhất, và nếu cách ấy chỉ đứng vững
nhờ một việc kiểm chứng đã làm ở cách đắt hơn, thì phải có ở đâu đó một khối
thí nghiệm có người làm nền cho toàn bộ Phần~IV. Câu hỏi thực nghiệm là khối
ấy có tồn tại không, và nó lớn tới đâu.

\index{lời giải thích!hai nhánh đánh giá của nó}%
Muốn hỏi được câu ấy thì phải tách rời hai nhánh mà một lời giải thích có thể
đi qua, vì hai nhánh này thường được gọi chung là \enquote{đánh giá} rồi bị
lẫn vào nhau. Hình~\ref{fig:ch12-hai-nhanh-danh-gia} vẽ chúng. Nhánh thứ nhất
là nhánh tự động: lời giải thích được đưa vào một thủ tục tính và ra một điểm
số, không ai được hỏi ý kiến, và đây là nhánh mà
chương~\ref{ch:do-do-trung-thuc} tới chương~\ref{ch:tro-choi-attribution} đã
đi. Nhánh thứ hai là nhánh người: lời giải thích được đưa cho một người đọc và
ra một đánh giá cảm nhận, chẳng hạn người ấy thấy nó dễ hiểu hay không, thấy
nó hợp lý hay không.

\begin{figure}[htbp]
  \centering
  \begin{tikzpicture}[
  every node/.style={font=\footnotesize, align=center},
  box/.style={draw=black!65, rounded corners=2pt, fill=black!5,
    minimum width=24mm, minimum height=11mm},
  outbox/.style={box, minimum width=21mm, fill=black!12},
  flow/.style={-{Stealth[length=2mm]}, thick, black!70},
  missing/.style={{Stealth[length=2mm]}-{Stealth[length=2mm]}, thick,
    black!70, densely dotted},
  gap/.style={-{Stealth[length=2mm]}, thick, black!70, densely dotted},
  lab/.style={font=\scriptsize, align=center, black!55}
]
  \node[box] (model) at (0,1.7) {Mô hình $f$};
  \node[box] (expl) at (0,-0.5) {Lời giải thích};

  \node[box, minimum width=25mm] (metric) at (4.3,1.0) {Chỉ số tự động};
  \node[box, minimum width=25mm] (human) at (4.3,-2.0) {Người đánh giá};

  \node[outbox] (score) at (8.2,1.0) {Điểm số};
  \node[outbox] (rating) at (8.2,-2.0) {Đánh giá\\cảm nhận};

  \draw[flow] (model) -- (expl);
  \draw[flow] (expl) to[out=60, in=180] (metric);
  \draw[flow] (expl) to[out=-60, in=180] (human);
  \draw[flow] (metric) -- (score);
  \draw[flow] (human) -- (rating);

  \draw[missing] (score) -- (rating);
  \node[lab, text width=23mm, anchor=east] at (7.0,-0.5)
    {phép đối chiếu\\giữa hai nhánh};

  \draw[gap] (expl) to[out=170, in=190, looseness=1.7] (model);
  \node[lab, text width=17mm, anchor=east] at (-2.4,0.6)
    {độ trung thực\\(định nghĩa~\ref{def:ch01-do-trung-thuc})};
\end{tikzpicture}

  \caption{Nét chấm là thứ đang thiếu, theo quy ước mà
  hình~\ref{fig:ch10-vong-bien-the} đặt. Thiếu ở hai chỗ khác nhau: phép đối
  chiếu giữa hai nhánh, mà mục~\ref{sec:ch12-con-so-va-cai-no-dem} đếm tần
  suất và mục~\ref{sec:ch12-ket-qua-va-ranh-gioi} đọc kết quả, và quan hệ
  giữa lời giải thích với mô hình, thứ mà không nhánh nào trong hai nhánh đo
  trực tiếp. Bản gộp của sách từ hai bài báo của chương
  này~\autocite{p24humangap,p25userperception}.}
  \label{fig:ch12-hai-nhanh-danh-gia}
\end{figure}

Hai nhánh ấy trả lời hai câu hỏi khác nhau, và
chương~\ref{ch:giai-thich-de-lam-gi} đã tách sẵn hai câu hỏi đó bằng cặp khái
niệm độ trung thực và \tn{tính hợp lý}{plausibility}. Một điểm số ở nhánh trên
nói về quan hệ giữa lời giải thích và mô hình. Một đánh giá cảm nhận ở nhánh
dưới nói về quan hệ giữa lời giải thích và người đọc. Người đọc không nhìn
thấy bên trong mô hình, nên dù có hỏi bao nhiêu người thì câu trả lời của họ
cũng không trực tiếp cho biết lời giải thích có phản ánh đúng mô hình hay
không; đó là lý do mũi tên vòng lên phía mô hình trong
hình~\ref{fig:ch12-hai-nhanh-danh-gia} vẫn là nét chấm ngay cả khi cả hai
nhánh đều chạy.

Vậy thì hỏi người để làm gì, nếu người không phán được điều mà Phần~IV đang
truy? Câu trả lời nằm ở chỗ khác, và nó là chỗ mà quan hệ vay mượn ở đầu mục
này chỉ ra. Không ai đòi người đọc phán về độ trung thực. Cái nhánh dưới làm
được là kiểm chứng chính nhánh trên: nếu một chỉ số tự động được đề xuất như
đại lượng thay thế cho chất lượng của lời giải thích \emph{đối với người dùng},
thì phép thử tối thiểu cho lời đề xuất ấy là cho thấy điểm số của nó đi cùng
đánh giá của người trên cùng những lời giải thích. Phép thử ấy chính là mũi tên
chấm nằm ngang giữa hai nhánh trong hình.

Cần giữ đúng phạm vi của phép thử ấy, vì nó không phải điều kiện xác nhận cho
mọi chỉ số. Một chỉ số được đề xuất để đo độ trung thực đang phát biểu về quan
hệ giữa lời giải thích và mô hình, và mục~\ref{sec:ch01-do-trung-thuc-va-tinh-hop-ly}
đã tách quan hệ ấy khỏi quan hệ với người đọc. Đòi một chỉ số như thế phải tương
quan với cảm nhận là đòi nó khớp với tính hợp lý, tức đúng thứ mà cả Phần~IV
đang cẩn thận không đồng nhất với độ trung thực; một chỉ số trung thực hoàn hảo
vẫn có thể không tương quan với cảm nhận, vì người đọc không nhìn thấy bên trong
mô hình. Bằng chứng xác nhận cho loại chỉ số ấy là ground truth dựng sẵn hoặc
một can thiệp có thiết kế, tức đúng đường mà
chương~\ref{ch:chi-so-khong-do-do-trung-thuc} đã đi. Vậy nhánh dưới xác nhận
được lời đề xuất \enquote{chỉ số này thay cho chất lượng người dùng cảm nhận},
và không xác nhận được lời đề xuất \enquote{chỉ số này đo độ trung thực}. Chương
này đếm xem phép thử ấy đã được chạy bao nhiêu lần, và giữ nguyên ranh giới vừa
nói khi đọc con số đếm được. Mục~\ref{sec:ch12-mot-cuoc-kiem-dem} và
mục~\ref{sec:ch12-con-so-va-cai-no-dem} đếm xem nó đã được chạy bao nhiêu lần,
còn mục~\ref{sec:ch12-mot-phep-doi-chieu} và
mục~\ref{sec:ch12-ket-qua-va-ranh-gioi} đọc một lần nó được chạy.
